\documentclass[11pt]{article}

\usepackage[margin=2.5cm]{geometry}
\usepackage[T1]{fontenc}
\usepackage{amsmath,amssymb}
\usepackage{xcolor}
\usepackage{parskip}
\usepackage{hyperref}

\definecolor{accentc}{HTML}{2A4FC7}
\hypersetup{colorlinks=true, linkcolor=accentc, urlcolor=accentc,
  pdftitle={CDP Active Recall}, pdfauthor={RAI}}

\newcommand{\qa}[2]{\par\noindent\textbf{#1.} #2\par\smallskip}

\title{Active Recall --- Causal Dependence Plots}
\author{}
\date{\today}

\begin{document}
\maketitle

\paragraph{How to use this.} This is a separate, standalone file from
\texttt{main.tex} on purpose --- \texttt{main.tex} is the reference to
consult, this one is for testing yourself \emph{without} the reference in
front of you. For each question: cover the Answers section (a ruler, a
folded page, or just don't scroll there), say or write your answer from
memory, \emph{then} check. Getting one wrong is the useful part, not a
failure --- it tells you exactly which page of \texttt{main.pdf} to reread.
Revisit questions you missed sooner than ones you got right (that's the
whole idea behind spaced repetition). Rebuild with
\texttt{latexmk -pdf recall.tex} after adding questions of your own ---
and you should add your own as you keep reading; this seed set stops at
what's in the notes as of today.

\tableofcontents

\section*{Questions}
\addcontentsline{toc}{section}{Questions}

\subsection*{Notation \& background}
\qa{1}{What does $\operatorname{do}(A=a)$ mean, and how is it different
  from conditioning on $A=a$?}
\qa{2}{Graphically, what does $\operatorname{do}(A=a)$ do to the DAG?}
\qa{3}{In one sentence, what is confounding?}
\qa{4}{In the running example, name the two separate effects age has that
  together create confounding between exercise and blood pressure.}
\qa{5}{What are the two conditions a covariate set $W$ must satisfy to be
  ``valid for adjustment''?}
\qa{6}{What can go wrong if you adjust for a variable that's a descendant
  of $A$? Give both failure modes.}
\qa{7}{What can go wrong if a backdoor path is left unblocked by your
  chosen adjustment set?}
\qa{8}{Why does the set of \emph{all parents} of $A$ always satisfy the
  backdoor criterion, with no case-by-case graph analysis needed?}
\qa{9}{True or false: if two covariate sets are both valid for adjustment,
  they must give equally precise (equal-variance) estimates.}
\qa{10}{In the backdoor adjustment formula, why do we average using
  $P(W{=}w)$ (the marginal distribution) instead of $P(W{=}w\mid A{=}a)$?}
\qa{11}{Potential-outcomes language needs three assumptions. Name them, and
  say which one is ``the backdoor criterion, in different words.''}
\qa{12}{What is a ``non-descendant'' of $A$, and which later result cares
  about \emph{all} non-descendants specifically (not just the parents)?}

\subsection*{The big picture}
\qa{13}{What does the classic PDP silently assume about the other features
  when one feature is varied?}
\qa{14}{Name all four named dependence plots and, in one phrase each, what
  each one does to downstream features.}
\qa{15}{Which one of the four plots is mathematically identical to the
  classic PDP?}
\qa{16}{What is the ``constructed outcome'' $Y^*$, and why is treating the
  model's output this way useful?}
\qa{17}{In the extended-graph picture ($\hat Y$ added as a child of every
  feature), which two plots are node interventions, and which two are edge
  interventions?}
\qa{18}{Why don't NDDP/NIDP need cross-world assumptions, even though
  they're edge interventions on a node that's ``acting two ways at once''?}
\qa{19}{Does direct $+$ indirect always sum to total? Name the leftover
  term and say when it vanishes.}

\subsection*{Identification}
\qa{20}{What is the ``reduced form'' $Q(dr\mid a,w)$, and why does it
  matter that it's learnable from data without any structural equations?}
\qa{21}{Which three of the four plots need only the parents of $A$ plus the
  reduced form? Which one needs more, and what does it need instead?}
\qa{22}{What object from the mediation-analysis literature does the NIDP
  turn out to equal exactly?}
\qa{23}{For this term's paper, is the NIDP handled for continuous $A$ or
  discrete $A$ only?}

\subsection*{Estimation}
\qa{24}{Name the three classical estimators for one discrete level of $A$,
  and say which one is ``doubly robust.''}
\qa{25}{In the AIPW formula, what happens to a person's correction term if
  their actual $A_i \neq a$?}
\qa{26}{What does ``doubly robust'' mean, precisely --- right if which
  condition(s) hold?}
\qa{27}{What does cross-fitting let you do, and what does it \emph{not} do
  (``not magic'')?}
\qa{28}{Between the parents of $A$ and all non-descendants of $A$, which is
  guaranteed never to be less efficient in large samples (assuming $f$ uses
  every feature)?}
\qa{29}{Why doesn't adjusting for a larger covariate set make the
  \emph{population} positivity condition harder, even though it makes
  estimation harder in finite samples?}

\subsection*{Continuous \texorpdfstring{$A$}{A}}
\qa{30}{Why is the bin weight $g_k$ chosen \emph{in advance} instead of
  using the observed distribution of $A$ inside the bin?}
\qa{31}{What goes wrong with naive confidence bands around a smoothed curve
  $\theta_h$, and what's the fix?}
\qa{32}{Why can a smooth-curve estimator's assumptions fail for a model $f$
  that's a tree splitting on $A$?}

\subsection*{Propagation vs.\ sensitivity}
\qa{33}{Does a downstream confounder (one not involving $A$) invalidate the
  parent-adjustment formula for the \emph{total} curve? Why or why not?}
\qa{34}{Name one advantage propagation can have over adjustment when the
  ECM's parametric equations are correct.}
\qa{35}{Can an unmeasured confounder change the \emph{level} of the curve
  in a way sensitivity analysis needs to worry about? Why does the analysis
  focus on the slope instead?}

\subsection*{Practical rules \& scope}
\qa{36}{Why must explanatory curves be drawn on data $f$ never trained on?}
\qa{37}{Why does ``the reference population'' matter? Name two different
  populations a curve could be averaged over that would give different
  answers.}
\qa{38}{Name two things this term's paper explicitly does \emph{not}
  claim.}

\subsection*{Molnar's PDP chapter (IML book, ch.\ 19)}
\qa{39}{Write the partial dependence function $\hat f_S(x_S)$ as an
  expectation. What is being averaged over, and what is held fixed?}
\qa{40}{Give the Monte-Carlo estimator actually used in practice. Are the
  $x_C^{(i)}$ values in that formula made up, or are they real rows from the
  dataset?}
\qa{41}{In one sentence, how does a PDP curve relate to a bundle of ICE
  (individual conditional expectation) curves?}
\qa{42}{What is the ``independence assumption'' underlying a PDP, and why is
  it called the biggest limitation of the method?}
\qa{43}{Molnar's height/weight example: computing the PDP at height
  $=200$cm means averaging over the marginal distribution of weight. What
  concrete, unrealistic combination can this produce?}
\qa{44}{Name the alternative plot built specifically to fix the
  correlation problem, and in one phrase, what does it condition on instead
  of marginalizing?}
\qa{45}{How is the partial dependence value for one category of a
  categorical feature (e.g.\ ``summer'') computed?}
\qa{46}{Per Zhao \& Hastie (as Molnar cites them), in what sense is a PDP
  ``causal,'' and in what sense is it explicitly \emph{not}?}
\qa{47}{Give the classic PDP-cancellation example: a feature where half the
  data has a positive association and half has a negative one. What does
  the PDP curve look like, and what would you wrongly conclude from it?}
\qa{48}{ICE plots fix the problem in Q47 by showing heterogeneity, but they
  still can't show you one key thing. What is it?}
\qa{49}{Why is two features usually the practical ceiling for a PDP, even
  though higher-order PDPs are mathematically well-defined?}
\qa{50}{Which one of the four CDP-family plots (\S2, ``The big picture,''
  in \texttt{main.tex}) turns out to be mathematically identical to the
  classic PDP? Given your answer to Q42, what is the classic PDP
  \emph{implicitly} assuming about the rest of the DAG?}

\clearpage
\section*{Answers}
\addcontentsline{toc}{section}{Answers}

\subsection*{Notation \& background}
\qa{1}{$\operatorname{do}(A=a)$ \emph{sets} $A$ to $a$ by intervention ---
  as if you forced it. Conditioning on $A=a$ just looks at whoever
  happened to have that value already, which can be contaminated by
  whatever else predicts $A$.}
\qa{2}{It deletes every arrow \emph{into} $A$ (its value no longer depends
  on its natural causes), then lets everything downstream propagate
  forward as usual.}
\qa{3}{Confounding is when two things look related only because they share
  a common cause, not because one causes the other.}
\qa{4}{Age affects blood pressure directly (aging), and age separately
  affects how much people exercise --- since it drives both, it creates a
  spurious link between exercise and blood pressure on top of any real
  causal link.}
\qa{5}{(a) It contains no descendant of $A$; (b) it blocks every backdoor
  path from $A$ to the outcome.}
\qa{6}{If the descendant is a \emph{mediator} (on the causal path from $A$
  to the outcome), conditioning on it blocks part of the effect you're
  trying to measure. If it's a \emph{collider}, conditioning on it can
  manufacture a spurious association that wasn't there before.}
\qa{7}{Some confounding is left in the adjusted number --- it will still be
  biased, just as the naive (unadjusted) comparison was, only possibly by
  a smaller amount.}
\qa{8}{Any path into $A$ must pass through one of $A$'s parents
  immediately before reaching it, so conditioning on all of them blocks
  every such path automatically --- this follows from the graph structure
  alone, no extra reasoning needed per graph.}
\qa{9}{False. Valid means ``gives the correct number''; different valid
  sets can differ in variance (how precisely that correct number is
  estimated from finite data) --- that's a separate, ``efficiency''
  question.}
\qa{10}{Because $P(W{=}w\mid A{=}a)$ is the age mix \emph{among people who
  happen to have} $A=a$ --- exactly the thing confounding distorts. Using
  the marginal $P(W{=}w)$ forces every level of $A$ to be compared against
  the same, undistorted population mix.}
\qa{11}{Consistency, ignorability/unconfoundedness, and
  positivity/overlap. Ignorability ($\{R(a)\}\perp A\mid W$) is the
  backdoor criterion restated in potential-outcomes language.}
\qa{12}{Everything that isn't a descendant of $A$ (parents, and anything
  upstream of or unrelated to $A$). The covariate-selection result
  (\S4, ``Estimation,'' in \texttt{main.tex}) shows adjusting on
  \emph{all} non-descendants --- not just the parents --- is never worse
  for variance, when $f$ uses every feature.}

\subsection*{The big picture}
\qa{13}{That the other features stay fixed at their recorded values ---
  i.e.\ that nothing downstream would actually respond if the varied
  feature really changed.}
\qa{14}{TDP: let everything downstream respond. NDDP: hold downstream
  features at natural values. NIDP: keep your own value of $A$, but move
  downstream features to their counterfactual values. PCDP: set $A=a$ and
  also fix one specific downstream feature.}
\qa{15}{NDDP.}
\qa{16}{$Y^* = f(X)$: the model's own output, treated as if it were an
  ordinary outcome variable. It's useful because it lets you apply
  ordinary causal-inference machinery (backdoor adjustment, AIPW, ...) to
  explaining a model, not just to real-world outcomes.}
\qa{17}{Node interventions: TDP and PCDP. Edge interventions: NDDP and
  NIDP.}
\qa{18}{Because $\hat Y = f(X)$ is a \emph{known, deterministic} function
  --- the usual cross-world assumptions are needed to handle uncertainty
  about how a node would behave under two conflicting interventions
  simultaneously, but there's no such uncertainty when the node's behavior
  is a known formula.}
\qa{19}{No --- there's a leftover interaction term $J(a)$ that must be
  added to make the two sides balance. It vanishes only when $f$ is
  additive in $A$ and the downstream features (no interaction between
  them).}

\subsection*{Identification}
\qa{20}{The conditional law of the non-parent features given $(A,W)$. It
  matters because it can be estimated directly from data with ordinary
  regression-style tools --- you never have to fit the ECM's actual
  structural equations.}
\qa{21}{TDP, NDDP, and NIDP. PCDP needs more: the graph \emph{inside} the
  non-parent features $R$ (two different internal graphs consistent with
  the same data can give different PCDP answers).}
\qa{22}{The ``bridge mean'' $\Psi(a) = E[Y(A,Z(a))]$ from the
  population-intervention indirect-effect literature.}
\qa{23}{Discrete $A$ only --- the continuous-feature version is flagged as
  needing its own theory, not done this term.}

\subsection*{Estimation}
\qa{24}{Plain regression (regress $Y^*$ on $(A,Z)$, average); inverse
  propensity weighting; and AIPW, which combines both and is the doubly
  robust one.}
\qa{25}{It's switched off entirely --- the indicator $\mathbf 1(A_i{=}a)$
  is zero, so that person contributes only their regression guess
  $\hat\mu(a,Z_i)$, with no correction term.}
\qa{26}{The estimator is consistent if \emph{either} the outcome regression
  $\hat\mu$ or the propensity $\hat\pi$ is correctly specified --- it
  doesn't need both.}
\qa{27}{It lets you use flexible ML models for the nuisance functions
  ($\hat\mu,\hat\pi$) while still getting a valid confidence interval. It
  does \emph{not} remove the underlying accuracy requirement --- the two
  nuisance errors still need to shrink fast enough (their product must be
  $o_p(n^{-1/2})$); cross-fitting just removes a technical condition, not
  that requirement.}
\qa{28}{All non-descendants of $A$.}
\qa{29}{Because $P(a\mid\text{non-descendants}) = P(a\mid W)$ by the graph
  --- the extra variables don't change the population-level propensity at
  all. What gets harder is purely a finite-sample \emph{fitting} problem
  (a bigger regression is harder to estimate well), not a population-level
  identification problem.}

\subsection*{Continuous \texorpdfstring{$A$}{A}}
\qa{30}{Because weighting by the \emph{observed} distribution of $A$ inside
  the bin would make the target depend on the data-generating law itself,
  and would add an extra term to the standard error. A weight fixed in
  advance keeps the target a clean, regular, $\sqrt n$ quantity.}
\qa{31}{Any uncorrected smoother has bias about as large as its standard
  error at its best bandwidth, so naive bands \emph{under-cover} (an
  earlier prototype got 0.89 coverage instead of $\sim$0.95). Fix:
  debiasing, deliberate undersmoothing, or reporting bands for the smoothed
  $\theta_h$ itself rather than pretending they're bands for the
  unsmoothed $\theta$.}
\qa{32}{A tree that splits on $A$ can make the true curve $\theta$ itself
  jumpy/discontinuous in $a$, but the smooth-curve theory assumes $\theta$
  is actually smooth --- so the assumption the theory needs is violated.}

\subsection*{Propagation vs.\ sensitivity}
\qa{33}{No --- the parent-adjustment formula for the total curve is still
  correct whether or not that downstream confounder is drawn in the graph.
  It changes the data (and so the curve), but adjustment still recovers
  whatever the right curve is for that data.}
\qa{34}{If the ECM's equations are right, propagation can be more precise
  and can extrapolate outside the region where the data actually overlap
  (as declared model-based extrapolation).}
\qa{35}{No --- the \emph{level} depends on an arbitrary zero point, so it's
  not a meaningful target for sensitivity analysis. The analysis instead
  asks how strong a confounder would need to be to flip the \emph{slope},
  flatten the curve, or cross a threshold --- things that don't depend on
  an arbitrary reference point.}

\subsection*{Practical rules \& scope}
\qa{36}{Because a flexible model can nearly reproduce its own training
  outcomes (a forest is a classic case) --- evaluating $f$ on rows it
  trained on isn't a fair sample of what $f$ actually does on new data.}
\qa{37}{The outer average is over some distribution of $W$, and different
  populations (e.g.\ the training cohort vs.\ a held-out sample vs.\ a
  future deployment population) can have different distributions of $W$
  --- so the curve genuinely differs across them whenever the effect
  varies with $w$.}
\qa{38}{Any two of: that parent adjustment/AIPW/cross-fitting/dose-response
  bands are new ideas; that the method works without a causal model at
  all; that it's the first semiparametric inference for causal
  explainability; bands for individual per-person curves.}

\subsection*{Molnar's PDP chapter (IML book, ch.\ 19)}
\qa{39}{$\hat f_S(x_S) = E_{X_C}[\hat f(x_S, X_C)]$. You average over the
  distribution of the \emph{other} features $X_C$, while the feature(s) of
  interest $x_S$ are held fixed at the value you're plotting.}
\qa{40}{$\hat f_S(x_S) = \frac1n\sum_{i=1}^n \hat f(x_S, x_C^{(i)})$. The
  $x_C^{(i)}$ are real, observed rows from the dataset --- for each one you
  swap in the fixed $x_S$ and keep that row's actual other-feature values,
  then average the model's output across all $n$ rows.}
\qa{41}{A PDP is the average of all the individual ICE curves in the
  dataset --- one ICE curve per data point, PDP is their pointwise mean.}
\qa{42}{It's assumed that the feature(s) in $S$ are uncorrelated with the
  features in $C$. It's the biggest limitation because real features are
  usually correlated, and averaging over the marginal (unconditional)
  distribution of $C$ then mixes in combinations that don't actually occur
  together.}
\qa{43}{A person who is 200cm tall but weighs under 50kg --- the PDP at
  height$=200$ averages over \emph{all} observed weights, including ones
  that in reality never co-occur with that height.}
\qa{44}{Accumulated Local Effect (ALE) plots. They average over the
  \emph{conditional} distribution of $C$ given $S$ (using only nearby,
  realistic combinations) instead of the marginal distribution of $C$.}
\qa{45}{Force every instance in the dataset to have that category (e.g.\
  set every row's season to ``summer''), run the model on all of them, and
  average the predictions.}
\qa{46}{It's causal \emph{for the model}: you are literally intervening on
  a feature and reading off how the model's output responds, which is a
  genuine causal statement about $\hat f$ as a function. It is \emph{not}
  necessarily causal for the real world the model was trained to predict.}
\qa{47}{If one half of the data has a positive slope and the other half an
  equally strong negative slope, the PDP curve can come out flat (the two
  halves cancel in the average). You would wrongly conclude the feature has
  no effect on the prediction, when in fact it has a strong effect that
  differs in direction across subgroups.}
\qa{48}{Which \emph{direction} (sign) the effect has for which subgroup ---
  ICE curves show that the lines fan out (heterogeneity exists) but not, by
  themselves, which lines belong to ``positive-effect'' people vs.\
  ``negative-effect'' people, especially when that split depends on yet
  another feature.}
\qa{49}{One feature gives a 2D plot and two features a 3D plot, which is
  about the limit of what a static image (or a human) can read directly.
  Higher-order PDPs are computable as functions --- there's just no
  meaningful way left to draw or perceive them.}
\qa{50}{NDDP (see recall.tex \S``Big picture'' above, Q15). It tells you the
  classic PDP is implicitly assuming there's no arrow from the varied
  feature to anything downstream that would actually respond --- i.e.\ it's
  silently treating the intervention as one where every other feature stays
  at its natural/recorded value, exactly the NDDP construction, not a full
  ``let everything respond'' intervention (that would be the TDP).}

\end{document}
